\begin{abstract}
Wikipedia is the largest open knowledge corpus, widely used worldwide and serving as a key resource for training large language models (LLMs) and retrieval-augmented generation (RAG) systems. Ensuring its accuracy is therefore critical. But how accurate is Wikipedia, and how can we improve it?

We focus on \textit{inconsistencies}, a specific type of factual inaccuracy, and introduce the task of \emph{corpus-level inconsistency detection}. We present \system, an agentic system that combines LLM reasoning with retrieval to surface potentially inconsistent claims along with contextual evidence for human review. In a user study with experienced Wikipedia editors, 87.5\% reported higher confidence when using \system, and participants identified 64.7\% more inconsistencies in the same amount of time.

Combining \system with human annotation, we contribute \dataset, the first benchmark of real Wikipedia inconsistencies. Using random sampling with \system-assisted analysis, we find that at least 3.3\% of English Wikipedia facts contradict another fact, with inconsistencies propagating into 7.3\% of \feverous and 4.0\% of AmbigQA examples. Benchmarking strong baselines on this dataset reveals substantial headroom: the best fully automated system achieves an AUROC of only 75.1\%.

Our results show that contradictions are a measurable component of Wikipedia and that LLM-based systems like \system can provide a practical tool to help editors improve knowledge consistency at scale.~\footnote{Dataset and code are available at \url{https://github.com/stanford-oval/inconsistency-detection}.}
\end{abstract}
